\sectionguard\subsection{The Greek Poets and Philosophers}\label{sec:bc-greeks}

The Greeks spoke for many centuries of beings set between God and men.
Hesiod sang of spirits who watch over mortal men; Socrates obeyed a sign
that came to him from boyhood; Plato taught that the daemonic nature
carries men's prayers to the gods and the gods' commands to men, that
each soul has a guardian allotted to it, and that the Maker of the world
made lesser gods who are immortal by his will; Aristotle proved that the
heavens are moved by substances without matter; and the later Platonists
built from these a whole ladder of gods, angels, and daemons, with rites
to climb it. The Fathers knew these writings well, and many of them
had been schooled in them. They received what the philosophers had
glimpsed of the created spirits, corrected it by revelation, and
unmasked the gods of the nations as the fallen angels, for ``all the
gods of the Gentiles are devils'' (Ps 95[96]:5) and ``the things which
the heathens sacrifice, they sacrifice to devils and not to God'' (1~Cor
10:20). Scripture itself gives the pattern: at Athens Paul quoted the
Greek poets, ``For we are also his offspring'' (Acts 17:28), and among
those who believed was ``Dionysius the Areopagite'' (Acts 17:34), under
whose name the Church received the \emph{Celestial Hierarchy}
(\S\ref{sec:ch}).

Aquinas states the rule of this reception at the opening of his treatise
\emph{On Separate Substances}, written, he says, because he could not be
present at the solemn office of the angels:
\begin{quote}
\latin{Intendentes igitur sanctorum Angelorum excellentiam utcumque
depromere, incipiendum videtur ab his quae de Angelis antiquitus humana
coniectura aestimavit: ut si quid invenerimus fidei consonum,
accipiamus; quae vero doctrinae repugnant Catholicae, refutemus.}
(\emph{De substantiis separatis}, prooem.)
\end{quote}
Intending to set forth, as far as may be, the excellence of the holy
angels, he begins with what human conjecture thought of the angels in
ancient times, so that whatever agrees with the faith may be accepted and
whatever is repugnant to Catholic doctrine refuted. The Fathers had done
the same, witness by witness.

\subsubsection{Hesiod and the watchers over men}\label{sec:bc-hesiod}

The oldest Greek witness is Hesiod's \emph{Works and Days}. After the
golden race of men who lived under Cronos had passed from the earth,
the poet says,
\begin{quote}
they are called pure spirits dwelling on the earth, and are kindly,
delivering from harm, and guardians of mortal men; for they roam
everywhere over the earth, clothed in mist and keep watch on judgements
and cruel deeds, givers of wealth. (Hesiod, \emph{Works and Days}
121--126, tr.\ Evelyn-White)
\end{quote}
Later in the poem he warns unjust princes that ``the deathless gods are
near among men,'' and that ``upon the bounteous earth Zeus has thrice
ten thousand spirits, watchers of mortal men, and these keep watch on
judgements and deeds of wrong as they roam, clothed in mist, all over
the earth'' (252--255). Hesiod thus holds three things that the Fathers
will recognize: a great multitude of spirits, a charge over men, and a
watch kept over justice. Scripture uses the same word for an angel.
Daniel saw in his vision ``a watcher, and a holy one'' come down from
heaven, and the sentence against the king is ``the decree by the
sentence of the watchers, and the word and demand of the holy ones''
(Dan 4:10, 14); the Douay note calls the watcher ``A vigilant angel.''

Plato reads Hesiod in the \emph{Cratylus}. Asked what the name
``demon'' means, Socrates recalls the golden race and quotes the lines,
``They are holy demons upon the earth, Beneficent, averters of ills,
guardians of mortal men,'' and concludes: ``he called them demons,
because they were daemones (knowing or wise)~\dots\ every wise man who
happens to be a good man is more than human (daimonion) both in life and
death, and is rightly called a demon'' (\emph{Cratylus} 397e--398c,
tr.\ Jowett). For Plato, then, the daemon is a knower, and the good man
becomes one after death. Plutarch in turn names Hesiod the first to
distinguish four rational natures, gods, daemons, heroes, and men
(\emph{De defectu oraculorum} 10).

The Fathers took up Hesiod's lines and the etymology together.
Lactantius quotes the poet's verse on the spirits who dwell on the
earth as guardians of mortal men, \latin{Hesiodus ita tradit}, and
explains it: ``And this is said for this purpose, because God had sent them as guardians
to the human race; but they themselves also, though they are the
destroyers of men, yet wish themselves to appear as their guardians,
that they themselves may be worshipped, and God may not be worshipped''
(\emph{Divinae institutiones} II.14.7--8, Brandt; ANF II.15;
\latin{custodes eos humano generi deus miserat}). Lactantius sets the
sending of these guardians within the history of the Watchers
(\S\ref{sec:before-christ}) and reads the poet by it: Hesiod saw truly
that God set guardians over men, and the fallen spirits have taken the
guardians' name to steal the worship due to God. He also receives the
etymology: the grammarians say that daemons are named \latin{quasi}
\gk{da\=emonas}\latin{, id est peritos ac rerum scios}, skilled and
acquainted with matters; they know many future things, but not all,
since God's counsel is hidden from them (II.14.6).

Eusebius turns Hesiod's line against the gods it was used to honour.
If there were good gods and daemons, why did they suffer men to
propitiate evil daemons with human sacrifice? ``And shall `The thrice
ten thousand guardians of mankind,' I mean their shepherds and
preservers~\dots\ deliver up their dearest ones to their enemies and
foemen, fierce as wild beasts~\dots?'' (\emph{Praeparatio evangelica}
IV.17, tr.\ Gifford). The guardians of the poets did not guard; they
were of one will with the demons.

Augustine sets the etymology beside the Apostle. ``They are called
demons from a Greek word meaning knowledge,'' \latin{Daemones enim~\dots\
dicuntur (quoniam vocabulum Graecum est) ab scientia nominati}; and
``the apostle, speaking with the Holy Spirit, says, `Knowledge puffeth
up, but charity buildeth up.'~\dots\ The demons, then, have knowledge
without charity, and are thereby so inflated or proud, that they crave
those divine honors and religious services which they know to be due to
the true God'' (\emph{De civitate Dei} IX.20, tr.\ Dods; 1~Cor 8:1).
The Platonist's praise of the daemon as a knower becomes, in Augustine,
the diagnosis of the demon's pride. Isidore hands the etymology to the
Latin Middle Ages in the same words as Lactantius, and adds that the demons' knowledge
comes partly from keener sense, partly from the experience of a very
long life, and partly \latin{per Dei iussum angelica revelatione}, by
angelic revelation at God's command (\emph{Etymologiae}
VIII.11.15--16). Aquinas receives it from Augustine: \latin{Daemones enim in Graeco a scientia nominantur quae sine
caritate, secundum sententiam apostoli, per superbiam inflat}: demons are
named in Greek from knowledge, which without charity, as the Apostle says,
puffs up through pride (\emph{De substantiis separatis} 20). The chain runs from the
\emph{Cratylus} through Lactantius, Augustine, and Isidore to Aquinas;
Plato and the Fathers agree on the word and differ on the thing, since
for the Fathers the name \emph{demon} belongs to the fallen.

Hesiod's guardians have their true fulfilment in the guardian angels.
Aquinas determines that ``each man has an angel guardian appointed to
him,'' taking as \emph{sed contra} Jerome on Matt 18:10, ``Great is the
dignity of souls, for each one to have an angel deputed to guard it from
its birth'' (\ST{I q.~113 a.~2}; \S\ref{sec:q113}). Where Hesiod made
the guardians men of a golden age raised to a higher state, the
Church confesses that the angels are a distinct creation, spiritual
creatures made from nothing (Lateran IV, DS~800;
\S\ref{sec:faith}).

\subsubsection{Socrates' sign}\label{sec:bc-socrates}

At his trial Socrates told the Athenians why he had kept out of public
life: ``You have heard me speak at sundry times and in divers places of
an oracle or sign which comes to me, and is the divinity which Meletus
ridicules in the indictment. This sign, which is a kind of voice, first
began to come to me when I was a child; it always forbids but never
commands me to do anything which I am going to do'' (Plato,
\emph{Apology} 31c--d, tr.\ Jowett). After his condemnation he took its
silence as a verdict: ``the oracle made no sign of opposition, either
when I was leaving my house in the morning, or when I was on my way to
the court~\dots\ It is an intimation that what has happened to me is a
good, and that those of us who think that death is an evil are in
error'' (40a--c). The indictment itself had charged him with introducing
new divinities.

Plutarch gave this sign its most careful philosophical account. In the
\emph{Discourse concerning Socrates's Daemon}, Simmias explains that it
was no apparition and no audible voice, ``but more probably a
declaration of a Daemon, by which the very thing that it would declare
was immediately and without audible voice represented to his mind''; for
``the understanding of a superior nature and a more divine soul may
excite an inferior soul, touching it from without,'' and ``the
conceptions of the Daemons carry a light with them, and shine to those
that are able to perceive them, so that there is no need of words''
(\emph{De genio Socratis} 20, tr.\ Goodwin). The account resembles
the Thomist doctrine of angelic illumination, in which ``the inferior
intellect is strengthened by the action of the superior intellect'' and
intelligible species are proposed to it (\ST{I q.~111 a.~1};
\S\ref{sec:illumination}).

The Fathers who name the sign divide in judging it. Justin honours
Socrates as a witness of the Word against the demons. The evil demons,
he says, had shown themselves to men who, ``not knowing that these were
demons,~\dots\ called them gods''; ``And when Socrates endeavoured, by
true reason and examination, to bring these things to light, and deliver
men from the demons, then the demons themselves, by means of men who
rejoiced in iniquity, compassed his death, as an atheist and a profane
person, on the charge that `he was introducing new divinities;' and in
our case they display a similar activity'' (\emph{1 Apology} 5, ANF).
Socrates ``taught men to reject the wicked demons,'' and Christ ``was
partially known even by Socrates (for He was and is the Word who is in
every man)'' (\emph{2 Apology} 10). Justin speaks of Socrates' mission,
not of his sign; the Latin apologists speak of the sign, and name it a
demon. Tertullian: ``The philosophers
acknowledge there are demons; Socrates himself waiting on a demon's
will. Why not? since it is said an evil spirit attached itself specially
to him even from his childhood---turning his mind no doubt from what was
good''; and in the same place, ``Plato also admits the existence of
angels,'' \latin{Angelos quoque etiam Plato non negavit}
(\emph{Apologeticum} 22.1--2). In the treatise \emph{On the Soul} he
calls the demon that ``clave to him from his boyhood'' ``the very worst
teacher certainly, notwithstanding the high place assigned to it by
poets and philosophers'' (\emph{De anima} 1). Minucius Felix writes that
``Socrates knew it, who, at the nod and decision of a demon that was at
his side, either declined or undertook affairs'' (\emph{Octavius} 26),
and Lactantius that Socrates spoke of a demon ever about him, who had
clung to him from boyhood and by whose will his life was ruled
(\emph{Div.\ inst.} II.14.9). Clement of Alexandria names the sign
without condemning it: reading Plato's allotted guide as the guardian
angel, he adds, ``Perhaps also the demon of Socrates suggested to him
something similar'' (\emph{Stromata} V.14; below). Augustine presses Apuleius on the same
point: the Platonist ``asserts most distinctly, and proves at great
length, that it was not a god but a demon,'' and so ``he ought not to
have entitled it, Concerning the God, but Concerning the Demon of
Socrates'' (\emph{De civ.\ Dei} VIII.14, tr.\ J.~J. Smith).

The Fathers thus hold distinct positions, each attributed: Justin sees
in Socrates a man who by reason turned from the demons and glimpsed the
Word; Clement names his sign without censure; Tertullian, Minucius, and
Lactantius judge the spirit that attended him to have been one of the
demons. All agree that the demons are real, that they sought worship
under the gods' names, and that the one teacher who expels them is
Christ.

\subsubsection{Plato: the daemonic between God and man}\label{sec:bc-plato}

Plato's teaching on the daemons lies scattered through the dialogues,
and the Fathers read it closely. It gathers under four heads: the
daemonic as interpreter, the guardian allotted to each soul, the made
gods of the \emph{Timaeus}, and the daemons as shepherds of men and
allies in an immortal war.

\paragraph{The interpreters.} In the \emph{Symposium} Diotima tells
Socrates that Love is ``a great spirit (daimon), and like all spirits he
is intermediate between the divine and the mortal.'' His power is this:
``He interprets~\dots\ between gods and men, conveying and taking across
to the gods the prayers and sacrifices of men, and to men the commands
and replies of the gods; he is the mediator who spans the chasm which
divides them~\dots\ For God mingles not with man; but through Love all the intercourse and
converse of God with man, whether awake or asleep, is carried on''; and
``these spirits or intermediate powers are many and diverse''
(\emph{Symposium} 202d--203a, tr.\ Jowett). Plato here sees two truths
and one error. The truths are that there are many spirits between God
and man, and that they carry to men what God commands; the error is that
God does not mingle with man.

Minucius Felix names the dialogue. Plato ``without hesitation, tell[s]
of both angels and demons. And in his Symposium also, does not he
endeavour to explain the nature of demons? For he will have it to be a
substance between mortal and immortal---that is, mediate between body
and spirit'' (\emph{Octavius} 26), and Lactantius likewise
(\emph{Div.\ inst.} II.14.9). Origen answers the same doctrine as Celsus had taken it
up. Christians, he says, ``acknowledge that angels are `ministering
spirits,' '' who ``ascend, bearing the supplications of men, to the
purest of the heavenly places~\dots\ and~\dots\ come down from these,
conveying to each one, according to his deserts, something enjoined by
God. Having thus learned to call these beings `angels' from
their employments~\dots\ [we are not] commanded to honour and worship in
place of God those who minister to us.'' For ``every prayer, and supplication, and intercession, and
thanksgiving, is to be sent up to the Supreme God through the High
Priest, who is above all the angels, the living Word and God''
(\emph{Contra Celsum} V.4, ANF; Heb 1:14). Origen keeps the office Plato
described---the ascent of prayers and the descent of gifts---and gives
it to the angels, named from their office; the mediation of worship he
reserves to Christ, and the name \emph{demon} he reserves to the wicked,
since it ``is always applied to those wicked powers, freed from the
encumbrance of a grosser body, who lead men astray'' (V.5).

Augustine answers the maxim that God does not mingle with man, which
Apuleius had repeated from Plato. It led the Platonists to place demons
in the air ``that they may carry to the gods the prayers of men, to men
the answers of the gods: for Plato held, they say, that no god has
intercourse with man,'' \latin{quoniam nullus deus miscetur homini, quod
Platonem dixisse perhibent} (\emph{De civ.\ Dei} VIII.18, tr.\ J.~J.
Smith). He answers with irony: \latin{Praeclara igitur sanctitas Dei,
qui non miscetur homini supplicanti, et miscetur daemoni arroganti}, a
fine holiness in God, who does not mingle with the man who prays and
does mingle with the arrogant demon (VIII.20). His
answer is the Incarnation: ``Good angels, therefore, cannot mediate
between miserable mortals and blessed immortals, for they themselves also
are both blessed and immortal; but evil angels can mediate, because they
are immortal like the one party, miserable like the other. To these is
opposed the good Mediator'' (IX.15), who became mortal for a time and
remains blessed for ever. Augustine expounds these books
at length (\S\ref{sec:fl-augustine}). Aquinas keeps what was true in
Diotima's word: ``The angels in their own nature stand midway between
God and men'' (\ST{I q.~64 a.~4}). The Church teaches that the angels are ``servants and messengers of God''
and that Christ is the centre of the angelic world (CCC~329, 331;
\S\ref{sec:faith}).

\paragraph{The allotted guardian.} Three dialogues give each soul a
daemon of its own. In the \emph{Phaedo} Socrates tells that ``after
death, as they say, the genius of each individual, to whom he belonged in
life, leads him to a certain place in which the dead are gathered
together, whence after judgment has been given they pass into the world
below''; the impure soul ``is after many struggles and many sufferings
hardly and with violence carried away by her attendant genius''
(\emph{Phaedo} 107d--108b, tr.\ Jowett, who renders \gk{daim\=on} as
``genius''). In the vision of Er at the end of the \emph{Republic}, the
prophet of Lachesis proclaims to the souls about to be born: ``Your
genius will not be allotted to you, but you will choose your genius
\dots\ the responsibility is with the chooser---God is justified''
(\emph{Republic} X 617d--e); when the souls had chosen, Lachesis ``sent
with them the genius whom they had severally chosen, to be the guardian
of their lives and the fulfiller of the choice'' (620d--e). The
\emph{Timaeus} places the daemon within: ``God gave the sovereign part
of the human soul to be the divinity of each one'' (\emph{Timaeus} 90a).

Clement of Alexandria reads the tale of Er as a witness to the angels.
By the fiery men who seize the wicked in Er's vision Plato means ``the
angels, who seize and punish the wicked'' (Ps 103[104]:4); and
``indicating `the angels' as the Scripture says, `of the little ones,
and of the least, which see God,' and also the oversight reaching to us
exercised by the tutelary angels, he shrinks not from writing~\dots\ [that Lachesis] sends along with each one, as his guide in life, and the
joint accomplisher of his purposes, the demon which he has chosen''
(\emph{Stromata} V.14, ANF; Matt 18:10). Clement receives Plato's
allotted guide as a figure of the guardian angel.

Justin reads another sentence of the same speech: ``Plato, when he
says, `The blame is his who chooses, and God is blameless,' took this
from the prophet Moses,'' who tells that God said to the first man,
``Behold, before thy face are good and evil: choose the good''
(\emph{1 Apology} 44).

Origen corrects the doctrine of the allotted daemon by the Gospel.
Against Celsus, who would have men give first-fruits to the demons set
over the world, he sets the thousands who stand before God and ``deem
it their office and duty to attend to'' those who pray, and answers the
Greeks with the Lord's word: ``Let the learned Greeks say that the human
soul at its birth is placed under the charge of demons: Jesus has taught
us not to despise even the little ones in His Church, saying, `Their
angels do always behold the face of My Father which is in heaven.' And
the prophet says, `The angel of the Lord encampeth round about them that
fear Him, and delivereth them' '' (\emph{Contra Celsum} VIII.34; Matt
18:10; Ps 33[34]:8). Tertullian makes the same correction from the
other side. Every heathen birth is consecrated to idols, and ``On this
principle of early possession it was that Socrates, while yet a boy, was
found by the spirit of the demon. Thus, too, is it that to all persons
their genii are assigned, which is only another name for demons''
(\emph{De anima} 39). For Tertullian the Roman \latin{genius} is a demon
that seizes the child at birth; for Origen the good guardian given to
the little ones is an angel. Both deny the name \emph{demon} to any
guardian sent by God.

The \emph{Phaedo}'s guide of the dead has its parallel in the Gospel,
where Lazarus ``was carried by the angels into Abraham's bosom'' (Luke
16:22). Tertullian names the pagan figure: at death the soul rejoices or
fears ``as soon as it sees the very angel's face, that arraigner of
souls, the Mercury of the poets'' (\emph{De anima} 53); and where Rome
invoked goddesses of childbirth, ``We, on our part, believe the angels
to officiate herein for God'' (37). Aquinas determines what the Fathers
had gathered from the Gospel: each man has an angel guardian, appointed
``from the very moment of his birth,'' following Jerome against those
who placed it at baptism (\ST{I q.~113 aa.~2, 5};
\S\ref{sec:guardians}).

\paragraph{The made gods of the Timaeus.} In the \emph{Timaeus} the
Maker, having framed the visible gods of heaven, speaks of the rest:
``To know or tell the origin of the other divinities is beyond us, and
we must accept the traditions of the men of old time who affirm
themselves to be the offspring of the gods'' (\emph{Timaeus} 40d). Then
he addresses all the gods he has made: ``Gods, children of gods, who are
my works~\dots\ since ye are but creatures, ye are not altogether
immortal and indissoluble, but ye shall certainly not be dissolved
\dots\ having in my will a greater and mightier bond than those with
which ye were bound at the time of your birth'' (41a--b). Plato here teaches that the gods below
the Maker are made, and that their immortality rests on his will.

Athenagoras quotes the passage on ``the other demons'' in full and reads
Plato's reserve as a hidden judgment: ``because he thought it impossible
to believe that gods beget and are brought forth,~\dots\ on this account
it was that he declared it to be beyond his powers to know and to speak
concerning the origin of the other demons'' (\emph{Legatio} 23, ANF).
Eusebius hears irony in the same lines: Plato ``speaks thus against his
own judgement on account of the laws'' (\emph{Praep.\ ev.}\ XIII.1).

Augustine receives the made gods of the \emph{Timaeus} as the holy
angels, provided that they are held to be created and blessed by
cleaving to God: ``If the Platonists prefer to call these angels gods
rather than demons, and to reckon them with those whom Plato, their
founder and master, maintains were created by the supreme God, they are
welcome to do so, for I will not spend strength in fighting about words.
For if they say that these beings are immortal, and yet created by the
supreme God, blessed but by cleaving to their Creator and not by their
own power, they say what we say, whatever name they call these beings
by,'' \latin{hoc dicunt quod dicimus, quolibet eos nomine appellent}
(\emph{De civ.\ Dei} IX.23, tr.\ Dods). Scripture itself calls them gods, ``The God of gods, the
Lord hath spoken'' (Ps 49[50]:1), but ``the name of demon is so
detestable that we cannot bear in any sense to apply it to the holy
angels'' (IX.23).

Eusebius presses the correction further. The Hebrews teach that ``the
intermediate nature of rational beings is generated and not without
beginning,'' and that from the ``fall and transgression'' of spirits,
powers, angels, and archangels ``they derive the race of daemons.''
Plato, ``although, like the Hebrews, he supposes the rational natures
to be incorporeal and intelligible essences,~\dots\ does not mean to
admit that they have arisen out of nothing,'' but derives them from ``an
effluence of the First Cause.'' If the daemons flowed from the good
gods, whence came their wickedness? The Hebrews answer that ``having no
previous existence it has come into being by the effective power of the
Cause of all,'' and has not goodness ``inseparably attached'' to it as
God has (\emph{Praep.\ ev.}\ XIII.15). The
Church confesses the same in the profession of Lateran IV: God \latin{sua
omnipotenti virtute simul ab initio temporis utramque de nihilo condidit
creaturam, spiritualem et corporalem, angelicam videlicet et mundanam
\dots\ Diabolus enim et alii daemones a Deo quidem natura creati sunt
boni, sed ipsi per se facti sunt mali} (DS~800;
\S\ref{sec:faith}).

Aquinas reads the address of the \emph{Timaeus} otherwise than
Augustine. An objection to the incorruptibility of the angels quotes the address
and argues that ``gods such as these can only be understood to be the
angels''; Aquinas answers: ``By the expression `gods'
Plato understands the heavenly bodies, which he supposed to be made up of
elements, and therefore dissoluble of their own nature; yet they are for
ever preserved in existence by the Divine will'' (\ST{I q.~50 a.~5
obj.~2 and ad~2}). He determines that the angels are incorruptible by
their nature, being subsisting forms; the perfect immortality that
includes immutability they have by grace (a.~5 corpus and ad~1;
\S\ref{sec:q50}). Augustine takes the made gods of the \emph{Timaeus} to
be the angels; Aquinas takes them to be the heavenly bodies. Both hold
that whatever Plato meant, the angels are made by God and blessed in
him.

\paragraph{Shepherds and allies.} In the \emph{Laws} Plato tells that in
the age of Cronos, ``God, in His love of mankind, placed over us the
demons, who are a superior race, and they~\dots\ taking care of us and
giving us peace and reverence and order and justice never failing, made
the tribes of men happy and united''; Cronos did with men ``as we do with flocks of
sheep'' (\emph{Laws} IV 713c--e, tr.\ Jowett). The \emph{Statesman}
tells the same age: ``the several parts the universe were distributed
under the rule of certain inferior deities~\dots\ There were demigods,
who were the shepherds of the various species and herds of animals''
(\emph{Statesman} 271d). In the tenth book of the \emph{Laws}, arguing
against the impious, the Athenian asks whether one soul or more orders
the heavens: ``More than one~\dots\ we must not suppose that there are
less than two---one the author of good, and the other of evil''
(\emph{Laws} X 896e); and later: ``there is, as we affirm, an immortal
conflict going on among us, which requires marvellous watchfulness; and
in that conflict the Gods and demigods are our allies, and we are their
property'' (906a).

Clement and Eusebius both hear Scripture in the tenth book. For Clement,
``that the devil so spoken of by the Barbarian philosophy, the prince of
the demons, is a wicked spirit, Plato asserts in the tenth book of the
Laws,'' and in the lines on the immortal conflict Plato ``expressly
emits that apostolic sentiment, `Our contest is not with flesh and
blood, but principalities, with powers' '' (\emph{Stromata} V.14; Eph
6:12). Eusebius quotes the same passages: ``thousands of years before
Plato was born this doctrine also had been acknowledged by the
Hebrews,'' who call ``the adverse power devil, and the good powers
angels of God'' (Job 1:6); and Plato's phrase that men are ``the
possessions of gods and daemons'' paraphrases the Septuagint's Deut
32:8, ``according to the number of the angels of God'' (\emph{Praep.\
ev.}\ XI.26; \S\ref{sec:before-christ}).

Plato's picture of the ages under Cronos, with lesser powers set over
the parts of the world and the kinds of living things, has its
Christian counterpart in the angels' government of bodies. Aquinas
declares the continuity in plain words: ``The holy doctors held with the
Platonists that different spiritual substances were placed over
corporeal things.'' For this he cites Augustine, ``Every visible thing
in this world has an angelic power placed over it,'' and Origen on the
ass of Balaam, that ``the world has need of angels who preside over
beasts, and over the birth of animals, and trees, and plants, and over
the increase of all other things'' (\ST{I q.~110 a.~1 ad~3}, citing
\emph{De div.\ qq.~83} q.~79 and \emph{Hom.\ in Num.}\ 14;
\S\ref{sec:q110}).
Athenagoras had given the same office its Christian form: ``this is the
office of the angels,---to exercise providence for God over the things
created and ordered by Him; so that God may have the universal and
general providence of the whole, while the particular parts are provided
for by the angels appointed over them'' (\emph{Legatio} 24).

Aquinas marks the limit of this continuity. Plato, as Aquinas reads
him through Nemesius (cited under Gregory of Nyssa's name), divided
providence into three, and the third, ``over human affairs, he assigned
to demons, whom the Platonic philosophers placed between us and the
gods'' (\ST{I q.~22 a.~3}). Aquinas rejects the division: ``Plato's
opinion is to be rejected, because he held that God did not govern all
things immediately, even in the design of government'' (\ST{I q.~103
a.~6 ad~1}). God's providence, as the plan of order, reaches every
thing immediately; its execution he carries out through intermediaries,
``not on account of any defect in His power, but by reason of the
abundance of His goodness; so that the dignity of causality is imparted
even to creatures'' (\ST{I q.~22 a.~3}). Plato saw the ministers; he did
not see that the King governs every least thing himself.

\subsubsection{Aristotle and the movers of the heavens}\label{sec:bc-aristotle}

Aristotle reached the separate substances by another road. In the
twelfth book of the \emph{Metaphysics} he argues that the first mover,
unmoved and eternal, produces the first and single movement of the
heavens; but since the planets have other eternal movements, ``each of
these movements also must be caused by a substance both unmovable in
itself and eternal~\dots\ Evidently, then, there must be substances
which are of the same number as the movements of the stars, and in their
nature eternal, and in themselves unmovable, and without magnitude''
(\emph{Metaphysics} XII.8, 1073a, tr.\ Ross). Following the astronomers
Eudoxus and Callippus he counts the spheres, fifty-five or forty-seven,
and concludes: ``Let this, then, be taken as the number of the spheres,
so that the unmovable substances and principles also may probably be
taken as just so many; the assertion of necessity must be left to more
powerful thinkers'' (1074a). He ends with the tradition of his
ancestors:
\begin{quote}
Our forefathers in the most remote ages have handed down to their
posterity a tradition, in the form of a myth, that these bodies are
gods, and that the divine encloses the whole of nature. The rest of the
tradition has been added later in mythical form with a view to the
persuasion of the multitude and to its legal and utilitarian
expediency~\dots\ But if one were to separate the first point from these
additions and take it alone---that they thought the first substances to
be gods, one must regard this as an inspired utterance. (XII.8, 1074b)
\end{quote}
In the treatise \emph{On the Heavens} he adds that ``all men have some
conception of the nature of the gods, and all who believe in the
existence of gods at all, whether barbarian or Greek, agree in allotting
the highest place to the deity'' (\emph{De caelo} I.3, 270b, tr.\
Stocks), and that the stars must be conceived ``as enjoying life and
action'' (II.12, 292a).

Aristotle's separate substances enter the doctrine of the angels through
Aquinas, who reads the twelfth book with care and corrects it at three
points. First, the number. Aristotle ``strove to find out the number of
the separate substances according to the number of the first
movements''; Maimonides tried to harmonize this with Scripture by
holding that the angels, as immaterial substances, are as many as the
heavenly movements. Aquinas determines against both that ``the angels,
even inasmuch as they are immaterial substances, exist in exceeding
great number, far beyond all material multitude,'' on the authority of
Daniel, ``Thousands of thousands ministered to Him, and ten thousand
times a hundred thousand stood before Him'' (Dan 7:10), and of
Dionysius (\emph{CH} 14). Aristotle's own argument, he notes, would
conclude only ``if the separate substances were made for corporeal
substances''; ``Hence Aristotle says (Metaph. xii, text 44) that this
is not a necessary argument, but a probable one'' (\ST{I q.~50 a.~3
and ad~3}; \S\ref{sec:q50}). The \emph{Contra Gentiles} makes the same
correction and adds reasons of its own: the order of the universe asks
that the nobler things exceed the less noble in number, as the heavenly
bodies exceed the elements in magnitude; \latin{excedunt igitur in numero
intellectuales substantiae separatae omnium rerum materialium
multitudinem}, the separate intellectual substances therefore exceed in
number the multitude of all material things (\emph{SCG} II.92 n.~7).

Second, the range of their government. ``Aristotle~\dots\ laid down
(Metaph. xi, 8) that the heavenly bodies are moved by spiritual
substances~\dots\ But he did not say that there were any spiritual
substances with immediate rule over the inferior bodies, except perhaps
human souls''; because ``many things are done in the inferior bodies
besides the natural corporeal actions,'' the Christian doctrine holds
that the angels preside over the lower bodies as well (\ST{I q.~110
a.~1 ad~2}). Third, the demons. In \emph{On Separate Substances} Aquinas
weighs the two masters. Aristotle's way ``by motion'' is surer, since it
does not depart far from what the senses show, but it is less
sufficient than Plato's. Things appear \latin{in hominibus qui a
Daemonibus opprimuntur, et in magorum operibus}, in men oppressed by
demons and in the works of magicians, that only an intellectual
substance can do; some Aristotelians traced them to the stars, as
Porphyry's letter to Anebo shows, but their cause \latin{plane quis
poterit secundum Platonicos assignare, si dicantur haec per Daemones
procreari}, one can plainly assign with the Platonists, if they are said
to be produced by demons (\emph{De substantiis separatis} 2). He notes also that Aristotle's proof proceeds from the
eternity of motion, \latin{quae fidei veritati repugnat}, which is
repugnant to the truth of the faith, and shows that it stands without
it, since the uniformity of the heavenly motions equally requires a
mover of infinite power (ibid.). Of the demons, he adds,
\latin{nec ipse nec eius sequaces fecisse mentionem}, neither Aristotle
nor his followers is found to have made mention (4).

Aquinas finds three agreements between Plato and Aristotle on the
separate substances: in their mode of being, since both hold that all
below the first are one and good by participation in it; in the
condition of their nature, since both hold them free of matter yet
composed of potency and act; and in providence, since both hold that the
higher provide for the lower under one first good, \latin{sicut unum
imperatorem vel dominum}, as under one emperor or lord (\emph{De sub.\
sep.}\ 3). The first two agreements are the Thomist doctrine of the angelic
substance itself: ``Although there is no composition of matter and form
in an angel, yet there is act and potentiality'' (\ST{I q.~50 a.~2
ad~3}; \S\ref{sec:substance}). On Aristotle's ``inspired utterance'' his
commentary on the \emph{Metaphysics} observes that the ancients handed
on \latin{per modum fabulae} that the separate substances are gods;
\latin{Si autem solum primum principium vocetur Deus, est unus tantum
Deus}, if only the first principle is called God, there is one God
only; and if one takes from the fables only this, \latin{quod dii sunt
quaedam substantiae immateriales, putabitur esse dictum divine et
secundum verisimilitudinem}, that the gods are certain immaterial
substances, it will be thought to have been said divinely and with
likelihood (\emph{In Metaph.} XII lect.~10 n.~31).

\subsubsection{Plutarch and Apuleius: the daemons of the air}\label{sec:bc-middle-platonists}

The Platonists of the first and second centuries after Christ made
Plato's scattered sayings into a doctrine. Plutarch gives it in the
dialogue \emph{Why the Oracles Cease to Give Answers}.
Those ``who have ranked the genus of Daemons between that of Gods and
men have solved greater doubts and difficulties, as having found the
knot which does, as it were, join and hold together our society and
communication with them'' (\emph{De defectu oraculorum} 10, tr.\
Goodwin). Xenocrates, Plato's friend, likened the gods to the
equilateral triangle, men to the scalene, and daemons to the isosceles,
since they are ``endued with the passions and perturbations of the
mortal nature, and the force and power of the divine''; to deny the
daemons is to destroy ``that nature (as Plato says) which serves as an
interpreter and messenger between them both.'' God is not present in
oracles himself, but commits them ``to his officers, the Daemons, who
are the spies and scouts of the Gods, wandering and circuiting about at
their commands,---some beholding and ordering the sacred ceremonies and
oblations offered to the Gods, others being employed to revenge and
punish the high misdemeanors and enormous injustices of men'' (13). The
shameful stories told of the gods are ``passions fitting to be attributed not to
Gods, but to Daemons'' (15); and ``It is not only Empedocles who affirms
there are bad Daemons, but even Plato, Xenocrates, and Chrysippus''
(17). The oracles fail because ``the Daemons who are appointed for the
government and superintendency of oracles'' have departed (15). Then
Plutarch tells how a voice from the island of Paxi called to the pilot
Thamus, ``When you are arrived at Palodes, take care to make it known
that the great God Pan is dead,'' and how, when he cried it aloud,
many voices were heard lamenting; the story reached Tiberius (17).

Apuleius of Madaura set the same doctrine in Latin, in his treatise
\emph{On the God of Socrates}. He grants Plato's maxim, \latin{nullus
deus miscetur hominibus} (\emph{De deo Socratis} 4), and fills the gap
with the daemons:
\begin{quote}
Besides, there are certain divine powers of a middle nature, situate in
this interval of the air, between the highest ether and the earth
below, through whom our aspirations and our deserts are conveyed to the
Gods. These the Greeks call by name ``daemons,'' and, being placed as
messengers between the inhabitants of earth and those of heaven, they
carry from the one to the other, prayers and bounties, supplications and
assistance, being a kind of interpreters and message carriers for both.
(\emph{De deo Socratis} 6, tr.\ Bohn 1853)
\end{quote}
His definition became the classic one: ``demons are as to genus animated
beings, as to mind rational, as to feelings passive, as to body aerial,
as to duration eternal,'' \latin{daemones sunt genere animalia, ingenio
rationabilia, animo passiva, corpore aeria, tempore aeterna} (13). The
daemon who guarded Socrates belongs to a higher class: ``Plato is of
opinion that a peculiar demon is allotted to every man, to be a witness
and a guardian of his conduct in life, who, without being visible to any
one, is always present, and is an overseer not only of his actions, but
even of his thoughts. But when life is finished, and the soul has to
return to its judges, then the demon who has presided over it
immediately seizes, and leads it as his charge to judgment'' (16). In
Latin, he says, the daemon that is each man's own soul may be called the
\latin{Genius} (15).

The Fathers answered both men by name. Eusebius quotes Plutarch's
chapters at length to prove from Plutarch's own pages that the oracles were
``the abodes of evil daemons'' (\emph{Praep.\ ev.}\ V.3--5, 16--17), and
he reads the voice at Paxi by its date:
\begin{quote}
But it is important to observe the time at which he says that the death
of the daemon took place. For it was the time of Tiberius, in which our
Saviour, making His sojourn among men, is recorded to have been ridding
human life from daemons of every kind: so that there were some of them
now kneeling before Him and beseeching Him not to deliver them over to
the Tartarus that awaited them. (\emph{Praep.\ ev.}\ V.17, tr.\ Gifford)
\end{quote}
The silence of the oracles, which Plutarch explained by the departure of
the daemons, Eusebius explains by the coming of Christ: ``all the
oracles and divinations of daemons are dead'' (V.1). Where Plutarch had said that
the fables of the gods are the sufferings of daemons, Eusebius adds that
``the deeds of Giants and Titans celebrated in song among the Greeks''
may be the story of the giants of Gen 6 (V.4), a reading that belongs
with the Watchers (\S\ref{sec:before-christ}).

Augustine answers Apuleius through the eighth and ninth books of the
\emph{City of God}. He reports the threefold division of rational
animals into gods in heaven, men on earth, and demons in the air
(\emph{De civ.\ Dei} VIII.14), and takes the definition apart clause by
clause: of what worth is it to be ``aerial in body,~\dots\ since a soul
of any kind whatsoever is to be set above every body?'' And ``how much
less are they really worthy of divine honor,---those aerial animals who
are only rational that they may be capable of misery, passive that they
may be actually miserable, and eternal that it may be impossible for
them to end their misery!'' (VIII.16). The demons' place in the air is
no rank but a prison: they are ``spirits most eager to inflict harm
\dots\ swollen with pride, pale with envy, subtle in deceit; who dwell
indeed in this air as in a prison~\dots\ because, cast down from the
height of the higher heaven, they have been condemned to dwell in this
element as the just reward of irretrievable transgression'' (VIII.22). The good
spirits whom the Platonists call good demons Scripture calls angels:
``we, following Scripture, according to which we are Christians, have
learned that some of the angels are good, some bad, but never have we
read in Scripture of good demons'' (IX.19).

Aquinas receives Augustine's judgment and states his own. An objection
in the \emph{Summa} quotes Augustine's commentary on Genesis, where the
demons are called ``animals of the atmosphere''; Aquinas answers that ``Augustine speaks, not as asserting the
fact, but merely using the opinion of the Platonists, who maintained that
there are some aerial animals, which they termed demons,'' and
determines that the angels have no bodies naturally united to them
(\ST{I q.~51 a.~1 ad~1}; \S\ref{sec:q51}). Origen, he adds, held the
contrary, ``being deceived here as he was also in many other points, by
following the opinions of the ancient philosophers'' (ibid.). The
Platonists' division of the heavens, on the other hand, Aquinas receives
at one point. They named as gods the intellectual substances above the
moon and as demons those beneath it, some good and some bad; Aquinas
cites Damascene (\emph{De fide orth.}\ II) for the view that the highest
of those who sinned was set over the terrestrial order, and judges:
``Nor is this opinion to be rejected as contrary to faith;
because the whole corporeal creation is governed by God through the
angels'' (\ST{I q.~63 a.~7}; \S\ref{sec:q63}). The dwelling of the demons
in the dark air until the judgment Aquinas takes from Augustine
(\S\ref{sec:q64}).

Apuleius's allotted demon, witness of thoughts and leader of the soul to
judgment, has its truth in the guardian angel of the Gospel
(\S\ref{sec:bc-plato}). Its name did not pass over. Lactantius says of
the demons that \latin{sibi geniorum nomen adsumunt; sic enim latino
sermone daemonas interpretantur}, they take to themselves the name of
genii, for so the Latins translate ``demons'' (\emph{Div.\ inst.}
II.14.12); Tertullian writes that the genii assigned to all persons are
``only another name for demons'' (\emph{De anima} 39). The guardian whom God gives is called an angel.

\subsubsection{Plotinus, Porphyry, and Iamblichus}\label{sec:bc-neoplatonists}

Plotinus wrote a treatise \emph{Of Our Individual Guardian}, the fourth
of the third \emph{Ennead}. He takes Plato's allotted daemon inward: ``What then is
our guardian? It is one of the powers of our soul~\dots\ our guardian
is the power immediately superior to the one that we exercise, for it
presides over our life without itself being active.~\dots\ If then we
live on the plane of the sense-life, our guardian is reason; if we live
on the rational plane, our guardian will be the principal superior to
reason~\dots\ Plato truly said that `we choose our guardian'; for, by
the kind of life that we prefer, we choose the guardian that presides
over our life'' (\emph{Enneads} III.4.3, tr.\ Guthrie, who renders
\gk{daim\=on} as ``guardian''). In the treatise on love he keeps the
daemons' middle rank: ``The deities are impassible, while the Guardians,
though eternal, can experience passions; placed beneath the deities, but
next to us, they occupy the middle place between deities and men''
(III.5.6). Porphyry's \emph{Life} tells that an Egyptian priest, calling
up Plotinus's guardian in a chapel of Isis at Rome, saw instead ``a
divinity of an order superior to that of guardians'' (\emph{Life of
Plotinus} 10).

Plotinus's guardian is a power of the soul; the guardian of the Gospel
is a person, an angel who beholds the face of the Father (Matt 18:10)
and is appointed to the man from birth (\ST{I q.~113 aa.~2, 5}).
Augustine received from Plotinus
something greater. The higher spirits, the Platonists taught, ``have
the same source of happiness as ourselves,---a certain intelligible
light, which is their God, and is different from themselves, and
illumines them that they may be penetrated with light''; Plotinus
compares God to the sun and the soul to the moon that has its light from
the sun, and says that the heavenly spirits derive ``their blessed life,
and the light of truth'' from the same source as we do; in this he is
``agreeing with the gospel where we read,~\dots\ `That was the true
Light which lighteth every man that cometh into the world' ''
(\emph{De civ.\ Dei} X.2; John 1:9). Angels and men have one light,
which is God.

Porphyry gives a full account of the good and the evil daemons. In the
second book \emph{On Abstinence} he distinguishes the souls that govern
the regions under the moon ``conformably to reason,'' which ``are to be
considered as good daemons,'' presiding over ``certain animals, or
fruits~\dots\ showers of rain, moderate winds, serene weather''; among
these ``those transporters, as Plato calls them,
[in his Banquet] are to be enumerated, who announce the affairs of men to
the Gods, and the will of the Gods to men'' (\emph{De abstinentia}
II.38, tr.\ Taylor). The others, who do not rule the spirit they are
joined to, are ``malefic demons.'' Eusebius quotes these chapters in
Gifford's rendering to convict the pagan cult by its own philosopher:
the evil daemons ``wish to be gods, and the power who presides over them
wishes to be thought the supreme god. These are they who delight `in
libations and burnt-offerings,' by which very things the spiritual and
bodily element is nourished and fattened'' (\emph{Praep.\ ev.}\ IV.22,
quoting \emph{De abst.}\ II.38--42). Eusebius draws the conclusion that
not only the poets but ``also those of the philosophers who were thought
to be earnest about the gods'' had worshipped ``not gods but wicked
daemons'' (IV.22).

Origen had answered the same doctrine as Celsus put it. Christians
maintain ``with regard not only to the fruits of the earth, but to every
flowing stream and every breath of air,'' that these come to be ``only in
consequence of the agency and control of certain beings whom we may call
invisible husbandmen and guardians; but we deny that those invisible
agents are demons''; to the demons, if they have any share, belong
``famine, blasting of the vine and fruit trees, pestilence'' as
executioners of God's judgments (\emph{Contra Celsum} VIII.31). Hence
``they who partake of corn and wine, and the fruits of trees, of water
and of air, do not feed with demons, but rather do they feast with
divine angels, who are appointed for this purpose'' (VIII.32). What
Porphyry gave to good daemons, Origen gives to the angels, and the gifts
they bring are received with thanksgiving to God. Of the demons
Origen affirms: ``according to our belief, it is true of all demons,
that they were not demons originally, but they became so in departing
from the true way'' (VII.69; Ps 95[96]:5); with God's favour, ``we
have also the good-will of all angels and spirits who are friends of
God,'' who join ``their own prayers and intercessions'' to ours
(VIII.64).

Iamblichus answered Porphyry's questions in the work \emph{On the
Mysteries}, and in it the words \emph{angel} and
\emph{archangel} enter the Platonic ladder. Porphyry had asked ``by what
indication the presence of a God, or an angel, or an archangel, or a
daemon, or a certain archon, or a soul, may be known''; Iamblichus
answers that the appearances of angels ``are more simple than those
of daemons, but are subordinate to those of the Gods; those of
archangels approximate in a greater degree to divine causes'' (\emph{De
mysteriis} II.3, tr.\ Taylor). The theurgist sought these appearances by
rites. Augustine names what they were. Porphyry ``promises some kind of
purgation of the soul by the help of theurgy,'' and ``distinguishes
angels from demons, asserting that the habitation of the latter is in
the air, while the former dwell in the ether and empyrean''; but the
theurgists, like the magicians they despise, are ``the slaves of the
deceitful rites of the demons whom they invoke under the names of
angels,'' and the lovely appearances of ``angels or gods'' seen in their
rites are what the Apostle means when he speaks of ``Satan transforming
himself into an angel of light'' (\emph{De civ.\ Dei} X.9--10; 2~Cor
11:14). In the letter to Anebo, Augustine grants, Porphyry wrote in ``a
better tone'' and ``repudiates all demons'' as drawn by the fumes of
sacrifice (X.11).

Porphyry himself had said ``that there are some angels who visit earth,
and reveal divine truth to theurgists, and others who publish on earth
the things that belong to the Father, His height and depth.'' Augustine
receives the second class and draws the consequence: ``Can we believe,
then, that the angels whose office it is to declare the will of the
Father, wish us to be subject to any but Him whose will they declare?
And hence, even this Platonist himself judiciously observes that we
should rather imitate than invoke them'' (X.26). The holy angels refuse
the sacrifice that belongs to God alone, ``in allegiance to whom they
themselves find their blessedness,'' and aid men ``with sincere
kindliness, and usurp over us no dominion, but declare to us Him under
whose rule we are then fellow-subjects'' (ibid.). What the Platonists
lacked was the way: ``You see in a fashion, although at a distance,
although with filmy eye, the country in which we should abide; but the
way to it you know not'' (X.29). The way is the Incarnation.

Aquinas takes up the whole Platonic series when he gives the causes of
idolatry. The Platonists ``said that there is one supreme god, the cause
of all things. After him they placed certain spiritual substances
created by the supreme god. These they called `gods,' on account of
their having a share of the godhead; but we call them `angels.' After
these they placed the souls of the heavenly bodies, and beneath these the
demons which they stated to be certain animal denizens of the air''
(\ST{II-II q.~94 a.~1}). Of idolatry itself, besides the causes on
man's part, there was a ``completive'' cause ``on the part of the demons, who
offered themselves to be worshipped by men, by giving answers in the
idols, and doing things which to men seemed marvelous. Hence it is
written (Psalm 95:5): `All the gods of the Gentiles are devils' ''
(a.~4). In the sixth age, he adds, ``idolatry was banished by the
doctrine and power of Christ, who triumphed over the devil'' (a.~4
ad~2). Eusebius had read the voice at Paxi in the same light.

\subsubsection{Proclus and the Dionysian corpus}\label{sec:bc-proclus}

Proclus, whom Aquinas calls \latin{Platonicus}, set out the school's
teaching on first causes in two hundred and eleven propositions, and the
corpus under the name of
Dionysius the Areopagite speaks the same language of procession, of
triads, of the higher illuminating the lower, and of participation in
the first good (\S\ref{sec:ch}). Aquinas knew both and says so. The
\emph{Book of Causes} he recognized as drawn from Proclus. In Greek, he
writes, there is the book of Proclus the Platonist, \latin{continens
CCXI propositiones, qui intitulatur elementatio theologica}; in Arabic
there is the book the Latins call \emph{On Causes}, which is not found in
Greek at all and so seems to have been excerpted \latin{ab aliquo
philosophorum Arabum ex praedicto libro Procli}, by one of the Arab
philosophers from Proclus's book (\emph{Super librum De causis},
prooem.).

Of Dionysius he writes that the difficulty of his books comes first
from this, \latin{quia plerumque utitur stilo et modo loquendi quo
utebantur Platonici, qui apud modernos est inconsuetus}, that he mostly
uses the style and manner of speech of the Platonists, which is unusual
among the moderns. The Platonists posited separate forms of natural
things, a man apart from matter, a horse apart from matter; this
\latin{fidei non consonat nec veritati}. But they also posited one first
principle, the very essence of goodness and unity and being, \latin{quod
dicimus Deum}, and from which all other things are good or one or being by
derivation; and \latin{quantum ad id quod dicebant de primo rerum
principio, verissima est eorum opinio et fidei Christianae consona},
as to what they said of the first principle of things, their opinion is
most true and consonant with the Christian faith (\emph{In librum B.
Dionysii De divinis nominibus}, prooem.). Hence
Dionysius calls God the Good itself, the super-good, the goodness of
every good.

Aquinas makes the Dionysian doctrine his chief authority on the Catholic
teaching about the angels: after surveying Plato and Aristotle he turns
to what the Christian religion asserts on each point, using chiefly
the teaching of Dionysius, \latin{qui super alios ea quae ad spirituales
substantias pertinent excellentius tradidit}, who handed on the things
that concern the spiritual substances more excellently than the others
(\emph{De sub.\ sep.}\ 18). There he corrects the Platonists' chain of causes: all
spiritual substances, not only the highest, are produced immediately by
God, as Dionysius teaches (ibid.; chs.~9--11). One Platonic principle he
keeps in the \emph{Summa}'s treatise on the angels, that ``the nearer things are to the one
first principle, the smaller they are in number''; it holds, he says,
``as regards the number of orders, as six administer and three assist''
(\ST{I q.~112 a.~4 ad~2}; \S\ref{sec:government}).

\subsubsection{The Fathers' account of the continuity}\label{sec:bc-fathers-account}

The Fathers did not only use the philosophers; they explained why the
philosophers had seen what they saw. They gave three accounts, and held
them together.

The first is that the Greeks learned from Moses. Justin says of Plato's
account of the making of the world, ``it was from our teachers---we mean
the account given through the prophets---that Plato borrowed his
statement that God, having altered matter which was shapeless, made the
world'' (\emph{1 Apology} 59); ``It is not, then, that we hold the same
opinions as others, but that all speak in imitation of ours'' (60).
Clement devotes long chapters to showing that ``the philosophy of the
Hebrews will be demonstrated beyond all contradiction to be the most
ancient of all wisdom'' (\emph{Stromata} I.21) and to the Greeks'
borrowings from the Hebrews (V.14), and among the
things borrowed he counts the angels who punish, the guardian angels,
the devil, and the spiritual combat (above). Eusebius gathers the whole
argument in the \emph{Preparation for the Gospel}, where ``the
philosophy of Plato in very many points contains a translation, as it
were, of Moses and the sacred writings of the Hebrews into the Greek
language'' (\emph{Praep.\ ev.}\ XIII, preface), and he copies Clement's
chapter at length (XIII.13).

Augustine tempers the claim by chronology. Some thought that Plato had
heard Jeremiah in Egypt, ``which opinion I myself have expressed in
certain of my writings. But a careful calculation of dates~\dots\ shows
that Plato was born about a hundred years after the time in which
Jeremiah prophesied'' and died before the Seventy translated the
Scriptures into Greek. Yet Augustine still inclines to think that Plato
learned their contents ``through an interpreter,'' moved above all by
the name God gave Moses through the angel, ``I am who am,'' a truth
``which Plato zealously held, and most diligently commended''
(\emph{De civ.\ Dei} VIII.11; Exod 3:14).

The second account is that the Word sows seeds of truth in every man.
``Whatever either lawgivers or philosophers uttered well, they
elaborated by finding and contemplating some part of the Word. But since
they did not know the whole of the Word, which is Christ, they often
contradicted themselves'' (Justin, \emph{2 Apology} 10). ``For each man
spoke well in proportion to the share he had of the spermatic word
\dots\ Whatever things were rightly said among all men, are the property
of us Christians'' (13). For Clement philosophy was God's gift to the
Greeks for their time: ``philosophy was given to the Greeks directly and
primarily, till the Lord should call the Greeks. For this was a
schoolmaster to bring `the Hellenic mind,' as the law, the Hebrews, `to
Christ' '' (\emph{Stromata} I.5).

The third account concerns the angels themselves. Clement records the
opinion that philosophy came to men as a thing stolen: ``It was then
some power or angel that had learned something of the truth, but abode
not in it, that inspired and taught these things, not without the Lord's
knowledge~\dots\ but without His prohibition'' (\emph{Stromata} I.17).
Elsewhere he gives the gentler account, grounded in the angels of the
nations: ``regiments of angels are distributed over the nations and
cities. And, perchance, some are assigned to individuals'' (VI.17); and
the Lord ``also gave philosophy to the Greeks by means of the inferior
angels. For by an ancient and divine order the angels are distributed
among the nations'' (VII.2). Clement sets the two accounts side by
side: philosophy stolen by a power that did not abide in the truth, yet
turned by Providence to good (I.17), and philosophy given through the
angels set over the nations (VII.2), which are the Septuagint's reading
of Deut 32:8 (\S\ref{sec:before-christ}).

Eusebius sets the Hebrews' doctrine of the angels beside the Greeks'.
The Hebrews, he says, teach that after God and his Word there are ``the
intelligent and rational Powers'' surpassing man's power to describe
``both for multitude and for variety of form,'' ordered like the stars,
for ``one star differeth from another star in glory''; Scripture names
them ``angels, and archangels, and spirits, and divine powers, and
heavenly hosts, principalities, and powers, and thrones, and dominions''
(\emph{Praep.\ ev.}\ VII.15; 1~Cor 15:41). ``Such are the doctrines
received from the Hebrews, which we have preferred to the erroneous
polytheism and daemonism of the Greeks, knowing and duly honouring
divine powers as servants and ministers of God the universal King, but
confessing Him alone as God, and worshipping Him alone'' (ibid.). The
fallen, led by him whom Isaiah calls the day star, received an abode in
Tartarus; ``of this race a small fragment left on the earth and in the
sublunar air to exercise the athletes of piety, has become a joint cause
of the polytheistic delusion of mankind,'' and the Hebrews alone from
the earliest ages taught ``that `all the gods of the nations are
daemons' '' (VII.16; Isa 14:12--14; Ps 95[96]:5). Plato ``alone of all
the Greeks touched the threshold of truth,'' yet ``sinks with the common
people of Athens'' when he honours the gods of the city
(XIII.14); therefore ``we must abandon this philosopher~\dots\ but must
honour Moses, and the Hebrew oracles'' (ibid.).

Augustine gives the Platonists the first place among the philosophers
and says why. ``It is evident that none come nearer to us than the
Platonists'' (\emph{De civ.\ Dei} VIII.5); whoever taught that God ``is
both the maker of all created things, the light by which things are
known, and the good in reference to which things are to be done,'' ``we
prefer these to all other philosophers, and confess that they approach
nearest to us'' (VIII.9). What they said truly is the Church's by right,
as Israel took the gold of Egypt:
\begin{quote}
Moreover, if those who are called philosophers, and especially the
Platonists, have said aught that is true and in harmony with our faith,
we are not only not to shrink from it, but to claim it for our own use
from those who have unlawful possession of it.~\dots\ some truths in
regard even to the worship of the One God are found among them. Now
these are, so to speak, their gold and silver, which they did not create
themselves, but dug out of the mines of God's providence which are
everywhere scattered abroad, and are perversely and unlawfully
prostituting to the worship of devils. (\emph{De doctrina christiana}
II.40.60, tr.\ Shaw)
\end{quote}
The image fits the angels exactly. The gold is the philosophers' knowledge
of a multitude of created spirits between God and man, bearers of his
will, guardians of men, movers of the heavens, ordered in ranks, some
fallen into envy and pride. The idols are the worship paid to those
spirits under the names of gods, the rites that fed the demons with
blood and smoke, and the theurgy that sought the angels by incantation.
The Fathers carried the gold out and left the idols. Aquinas, reading
the same philosophers eight centuries later, gathers the reception into
a sentence: ``The holy doctors held with the Platonists that different
spiritual substances were placed over corporeal things'' (\ST{I q.~110
a.~1 ad~3}); and of the Platonists' gods, ``we call them `angels' ''
(\ST{II-II q.~94 a.~1}).

\subsubsection{The philosophers and their reception by locus}\label{sec:bc-concordance}

\begin{fourstable}{Philosopher and locus}{Teaching}{Patristic reception}{Disposition}
Hesiod, \emph{Works and Days} 121--126, 252--255 & The golden race
become pure spirits, guardians of mortal men; thirty thousand watchers
over justice & Lactantius quotes 122--123: God sent guardians, whose
name the fallen usurp (\emph{Div.\ inst.} II.14.7--8); Eusebius turns
the thirty thousand guardians against the gods (\emph{Praep.\ ev.}\
IV.17) & Received as to guardians sent by God; corrected as to their
origin; the fallen unmasked (Lactantius). Guardian angels:
\ST{I q.~113 a.~2}\\
Plato, \emph{Cratylus} 397e--398c & Daemon named from knowing; the good
man becomes a daemon & Lactantius II.14.6; Augustine, \emph{De civ.\
Dei} IX.20; Isidore, \emph{Etym.} VIII.11.15 & Etymology received;
knowledge without charity puffs up (Augustine, 1~Cor 8:1). Aquinas,
\emph{De sub.\ sep.}\ 20\\
Plato, \emph{Apology} 31c--d, 40a--c & Socrates' sign from childhood,
which forbids & Justin, \emph{1~Apol.}\ 5, \emph{2~Apol.}\ 10;
Tertullian, \emph{Apol.}\ 22, \emph{De anima} 1; Minucius,
\emph{Oct.}\ 26; Lactantius II.14.9; Clement, \emph{Strom.}\ V.14 &
Disputed among the Fathers: Socrates a foe of the demons (Justin); his
spirit an evil demon (Tertullian, Minucius, Lactantius); named without
censure as a source of Plato's allotted guide (Clement)\\
Plato, \emph{Symposium} 202d--203a & The daemonic interprets between
gods and men; ``God mingles not with man'' & Minucius, \emph{Oct.}\ 26;
Origen, \emph{C.\ Cels.}\ V.4--5; Augustine, \emph{De civ.\ Dei}
VIII.18, 20, IX.15 & Office received and given to the angels (Origen);
maxim rejected; Christ the one Mediator (Augustine)\\
Plato, \emph{Republic} X 617d--e, 620d--e; \emph{Phaedo} 107d--108b &
A daemon allotted to each soul, guardian of its life and guide after
death & Clement, \emph{Strom.}\ V.14; Origen, \emph{C.\ Cels.}\ VIII.34;
Tertullian, \emph{De anima} 39, 53; Justin, \emph{1~Apol.}\ 44 &
Received as a figure of the guardian angel (Clement); corrected by Matt
18:10 (Origen); the pagan \latin{genius} rejected as a demon
(Tertullian). \ST{I q.~113 a.~5}\\
Plato, \emph{Timaeus} 40d--41b & Made gods, immortal by the Maker's
will; the origin of the other daemons beyond telling & Athenagoras,
\emph{Leg.}\ 23; Augustine, \emph{De civ.\ Dei} IX.23; Eusebius,
\emph{Praep.\ ev.}\ XIII.1, 15 & Received as the angels if created and
blessed by cleaving to God (Augustine); emanation corrected to creation
from nothing (Eusebius; Lateran IV, DS~800). Aquinas reads the gods as
heavenly bodies (\ST{I q.~50 a.~5 ad~2})\\
Plato, \emph{Laws} IV 713c--e; \emph{Statesman} 271d & Daemons set by
God as shepherds over men and kinds of living things & Athenagoras,
\emph{Leg.}\ 24, the angels' particular providence (parallel) & Received: ``The
holy doctors held with the Platonists'' (\ST{I q.~110 a.~1 ad~3}); the
threefold providence rejected (\ST{I q.~103 a.~6 ad~1})\\
Plato, \emph{Laws} X 896e, 906a & A good and an evil soul in the
heavens; gods and daemons our allies in an immortal conflict & Clement,
\emph{Strom.}\ V.14; Eusebius, \emph{Praep.\ ev.}\ XI.26 & Received as
from Moses: the devil and the spiritual combat (Eph 6:12); Job 1:6\\
Aristotle, \emph{Metaphysics} XII.8, 1073a--1074b & Unmoved movers as
many as the heavenly motions, by probable argument; the first
substances are gods & Aquinas, \ST{I q.~50 a.~3}; \emph{SCG} II.92;
\emph{De sub.\ sep.}\ 2--4; \emph{In Metaph.}\ XII lect.~10 &
Corrected: the angels exceed all material multitude (Dan 7:10); rule
the lower bodies too (\ST{I q.~110 a.~1 ad~2})\\
Plutarch, \emph{De defectu oraculorum} 10--17 & Daemons midway,
officers of the gods, some evil; oracles fail when they depart; ``Great
Pan is dead'' & Eusebius, \emph{Praep.\ ev.}\ V.1, 3--5, 16--17 &
Rejected as to worship: the oracles were demons', silenced under
Tiberius by Christ\\
Apuleius, \emph{De deo Socratis} 4, 6, 13, 16 & Aerial, passive,
eternal daemons carry prayers; each man's allotted guardian &
Augustine, \emph{De civ.\ Dei} VIII.14--22, IX.19 & Rejected: the air
is the demons' prison; good spirits are angels, never ``good demons.''
Aerial bodies: \ST{I q.~51 a.~1 ad~1}\\
Plotinus, \emph{Enneads} III.4.3; III.5.6 & The guardian a higher power
of one's own soul; daemons passible, between gods and men & Augustine,
\emph{De civ.\ Dei} X.2 (on illumination) & Parallel on the guardian;
received on the one light of angels and men (John 1:9)\\
Porphyry, \emph{De abstinentia} II.38--42; letter to Anebo & Good
daemons over fruits and weather; evil daemons crave worship and
sacrifice & Eusebius, \emph{Praep.\ ev.}\ IV.22; Origen, \emph{C.\
Cels.}\ VIII.31--32 (the same doctrine in Celsus); Augustine, \emph{De civ.\ Dei} X.9--11, 26 & Good
powers received as angels (Origen); demon-worship rejected; imitate, do
not invoke, the angels (Augustine)\\
Iamblichus, \emph{De mysteriis} II.3 & Gods, archangels, angels,
daemons, archons, souls, known by apparitions in theurgy & Augustine,
\emph{De civ.\ Dei} X.9--10 & Rejected: demons invoked under the names
of angels (2~Cor 11:14)\\
Proclus, \emph{Elements of Theology} (through the \emph{Liber de
causis}) & The order of first causes in 211 propositions, as Aquinas
describes it & Aquinas,
\emph{In De causis} prooem.; \emph{In De div.\ nom.}\ prooem.;
\emph{De sub.\ sep.}\ 18 & Received as to the first principle;
separate forms rejected; all spirits created immediately by God\\
\end{fourstable}
